\scanpage{14}
128\hfill J.-P. \textsc{Serre}

\textbf{Proposition 7.} -- (a) \textit{Le commutant de $G$ dans $T$ est égal à $\operatorname{End}_A(F)$.}

(b) \textit{Le commutant de $G$ dans $V$ est égal à $\operatorname{End}_A(F)\otimes\mathbf Q_p$.}

(c) \textit{Le commutant de $\mathfrak g$ dans $V$ est égal à $\operatorname{End}(F)\otimes\mathbf Q_p$.}

L'assertion (a) est démontrée dans [9], \S\ 4, cor. 1 au th. 4; (b) résulte de (a) par produit tensoriel avec $\mathbf Q_p$. Enfin, (c) résulte de (b) et du fait que le commutant de $\mathfrak g$ dans $V$ est égal à la réunion des commutants des sous-groupes ouverts de $G$.

Nous pouvons maintenant appliquer le théorème 3. On en tire:

\textbf{Théorème 4.} -- \textit{Faisons sur $F$ les hypothèses suivantes:}

(a) \textit{$V$ est un $G$-module semi-simple.}

(b) $\operatorname{End}(F)=\mathbf Z_p$.

(c) \textit{Les dimensions $n_1$ et $n_0$ de $F$ et de son groupe dual $F'$ sont $\geq1$ et premières entre elles.}

\textit{On a alors $G_{\mathrm{alg}}=\mathbf{GL}_V$, $\mathfrak g=\operatorname{End}(V)$ et $G$ est un sous-groupe ouvert de $\operatorname{Aut}(V)$.}

En effet, (H*.1) résulte de (a) et (b), compte tenu de la prop. 7, et (H*.2) et (H*.3) résultent de la prop. 6 et de (c).

\textbf{Remarque.} Les hypothèses (a), (b), (c) entraînent que $F$ est connexe (c'est donc un \textit{groupe formel}). En effet, si $F^0$ désigne la composante neutre de $F$, $V(F^0)$ est un sous-espace de $V(F)=V$ stable par $G$. Vu (a) et (b), on a donc $V(F^0)=0$ ou $V(F^0)=V$. Le premier cas entraîne $F^0=0$, d'où $n_1=0$, ce qui est exclu d'après (c). Le second cas entraîne $F=F^0$. De même, le dual de $F$ est un groupe formel.

Nous allons maintenant appliquer le théorème 4 aux \textit{groupes formels de dimension 1} (au sens de \textsc{Lubin} [2], [3]), i.e. au cas $n_1=1$. On a tout d'abord:

\textbf{Proposition 8.} -- \textit{Supposons $F$ connexe et de dimension 1. Pour tout $z\in V$, $z\ne0$, on a $\mathfrak g\cdot z=V$. En particulier, $V$ est un $\mathfrak g$-module simple.}

Par hypothèse, l'algèbre affine de $F$ est isomorphe à l'algèbre des séries formelles $A[[X]]$. Choisissons un tel isomorphisme (ce qui revient à considérer $F$ comme un groupe formel). La multiplication par $p$ dans $F$ est donnée par une série formelle
\[
f(X)=\sum_{n=1}^{\infty}a_nX^n,\qquad a_n\in A,\quad a_1=p.
\]
Puisque $F$ est de hauteur $h$, on a
\[
a_n\in\mathfrak m\quad\text{pour }n<p^h,\quad\text{et}\quad a_n\notin\mathfrak m\quad\text{pour }n=p^h,
\]
cf. \textsc{Lubin} [2].
\margin{13}
